% A resource theory of quantum energy teleportation.
% Build: latexmk -pdf main.tex   (figures are drawn by pgfplots from ./data/*.dat,
% which export_figdata.py writes from the candidate's archived ../../data/).
%
% BEFORE SUBMISSION (see ../../MANIFEST.md and the commit that added this file):
%  1. Read Fan et al. 2024 (PRA 110, 052424; QIP 23, 367) and Wu et al. 2024
%     (PRA 109, 062208) in full. They are journal-only (open-access retry 2026-09-23:
%     Unpaywall, Semantic Scholar, arXiv, 7 citing papers -- still closed, no preprint);
%     the Related-work paragraph describes them from their abstracts and from how
%     Refs. [Xie2025, Yan2026] cite them, and disclaims the two-qubit cases of
%     Propositions 1 and 3. Check in particular whether their "QET passive" states coincide with
%     the free states here and whether their correlation-matrix analysis already
%     contains Proposition 1 or Proposition 3 for the minimal model.
%  2. Have a quantum-energy-teleportation researcher read the whole draft.
%  3. The AI-use statement in the acknowledgments (set 2026-09-23) is the author's to
%     keep or edit; name any other AI tools that were used.
\documentclass[aps,pra,twocolumn,reprint,superscriptaddress,nofootinbib,floatfix]{revtex4-2}

\usepackage{amsmath,amssymb,amsthm}
\usepackage{bm}
\usepackage{graphicx}
\usepackage[colorlinks=true,linkcolor=blue!50!black,citecolor=blue!50!black,urlcolor=blue!50!black]{hyperref}
\usepackage{pgfplots}
\pgfplotsset{compat=1.18}
\usepgfplotslibrary{groupplots}

\theoremstyle{plain}
\newtheorem{theorem}{Theorem}
\newtheorem{proposition}{Proposition}
\newtheorem{corollary}{Corollary}
\theoremstyle{remark}
\newtheorem{remark}{Remark}

\newcommand{\Tr}{\operatorname{Tr}}
\newcommand{\id}{\mathrm{id}}
\newcommand{\Wto}{W_{\to}}
\newcommand{\WCK}{W_{\to}^{\mathrm{CK}}}
\newcommand{\Wloc}{W_{\mathrm{loc}}}
\newcommand{\HBl}{H_B^{\mathrm{loc}}}
\newcommand{\Abar}{\bar{\mathcal{A}}}
\newcommand{\E}{\mathcal{E}}
\newcommand{\lmin}{\lambda_{\min}}
\newcommand{\Hmin}{H_{\min}}
\newcommand{\ket}[1]{|#1\rangle}
\newcommand{\bra}[1]{\langle #1|}
\newcommand{\proj}[1]{|#1\rangle\langle #1|}
\newcommand{\sgn}{\operatorname{sgn}}
\newcommand{\spread}{\operatorname{spread}}

\begin{document}

\title{A resource theory of quantum energy teleportation}

\author{Matthew Ruckman}
\email{mruckman1@gmail.com}
\thanks{ORCID: \href{https://orcid.org/0009-0002-1723-3823}{0009-0002-1723-3823}}
\affiliation{Unaffiliated}

\date{\today}

\begin{abstract}
In quantum energy teleportation (QET) a receiver, Bob, extracts energy from a strongly
locally passive state after a sender, Alice, measures her part and sends him the outcome.
We formulate QET as a resource theory. The free operations are Alice's instruments whose
average cannot change Bob's energy (energy-non-signalling instruments), one-way classical
communication, and arbitrary channels on Bob's side, with the energy that Bob's device
receives booked. The QET potential $\Wto$, the most energy Bob can receive by free
operations, is a ledger monotone: no free operation increases $\Wto$ plus the energy
already delivered. When Alice's bond operators commute, the value of $\Wto$ over
commuting-Kraus instruments has a closed form, which recovers Hotta's optimum for the
minimal model. We prove that commuting-Kraus instruments exhaust the free class exactly
when the joint eigenvalue vectors of the bond operators are in convex position. In that
case $\Wto$ is linear in the state, the ledger holds with equality, and two-way classical
communication delivers no more energy to Bob than one-way communication. Off convex
position the closed form is not a monotone: in a star-coupled model, a free instrument
outside the commuting class raises it 6.15-fold. The QET surplus $\Wto-\Wloc$ is not a
monotone in either direction. For a single-qubit receiver we apply the local-ergotropy
formula of Salvia, De Palma, and Giovannetti branch by branch; entropic bounds on Bob's
yield are valid but loose, and we identify why. Numerical checks cover 12\,500 random free
operations and transverse-field Ising chains of up to 12 sites.
\end{abstract}

\maketitle

% ---------------------------------------------------------------------------------------
\section{Introduction}
\label{sec:intro}

Ground states of interacting many-body systems are, in general, strongly locally passive:
no operation on a subsystem lowers the total energy~\cite{Frey2014,Alhambra2019}. Hotta
observed that this energy is nevertheless available to local operations aided by classical
communication~\cite{Hotta2008,Hotta2009,Hotta2011review}. In quantum energy teleportation
(QET), Alice measures a local observable and sends the outcome to Bob, who applies an
outcome-dependent operation that lowers the energy near him. Alice's measurement injects an
energy $E_A$ near her, Bob extracts $E_B<E_A$ near him, and no energy travels between them
faster than the classical message. QET has been demonstrated in nuclear magnetic
resonance~\cite{RodriguezBriones2023} and on superconducting hardware~\cite{Ikeda2023};
Ref.~\cite{Ragula2025} reviews its applications.

Recent work varies the ingredients: the initial state~\cite{Wu2024,Yan2026}, Alice's
measurement~\cite{Xie2025}, and the roles of entanglement, coherence, steering, and
information teleportation~\cite{Fan2024steering,Fan2024resources,Wang2024}. Comparing such
variants requires fixing which operations are allowed and what is being spent. Resource
theories supply that language: a class of free operations, the states they cannot improve,
and monotones that quantify what the free operations consume~\cite{ChitambarGour2019}. This
paper builds such a theory for QET.

The free operations (Sec.~\ref{sec:setting}) are Alice's energy-non-signalling
instruments, one-way classical communication from Alice to Bob, and arbitrary channels on
Bob's side. An instrument is energy-non-signalling if its average leaves Bob's energy
unchanged on every state, the energetic form of no-signalling. The energy each of Bob's
channels delivers to his device is booked. The monotone is the QET potential $\Wto$, the
most energy Bob can receive by free operations. Our results are the following.
\begin{enumerate}\setlength{\itemsep}{1pt}\setlength{\parskip}{0pt}
\item $\Wto$ is a ledger monotone: no free operation increases $\Wto$ plus the energy
already delivered to Bob's device (Theorem~\ref{thm:ledger}). This is close to built in;
the content lies in the next items.
\item For commuting bond operators, the potential over commuting-Kraus instruments,
$\WCK$, is attained by Alice's projective bond measurement followed by Bob's reset to the
ground state of the conditional Hamiltonian that Xie, Sajjan, and Kais call his local
effective Hamiltonian~\cite{Xie2025} (Proposition~\ref{prop:closed}). We state it because
the next two items rest on it; for the minimal model it reproduces Hotta's optimum.
\item Commuting-Kraus instruments are the whole free class if and only if the joint
eigenvalue vectors of the bond operators are in convex position
(Theorem~\ref{thm:freeclass}).
\item $\WCK$ is a linear functional of the state. In convex position, therefore, the
ledger is an equality, and two-way classical communication delivers no more to Bob than
one-way communication (Theorem~\ref{thm:linear}, Corollary~\ref{cor:twoway}). Off convex
position the closed form is not a monotone (Sec.~\ref{sec:offconvex}).
\item The QET surplus $\Wto-\Wloc$ is not a monotone under the free operations, in either
direction, nor under local operations and classical communication (LOCC)
(Proposition~\ref{prop:surplus}).
\item For a single-qubit Bob, the local-ergotropy formula of Salvia, De Palma, and
Giovannetti~\cite{Salvia2023} gives his optimal unitary yield on each measurement branch.
We use it to show that the conditional rotation of the spin-chain protocol is optimal among
unitaries, and we show why entropic bounds on the yield are loose (Sec.~\ref{sec:single}).
\end{enumerate}

Everything is finite-dimensional, apart from a Gaussian analogue in
Sec.~\ref{sec:gaussian}. Nothing here concerns quantum fields or a thermodynamic limit.
The numerical checks (Sec.~\ref{sec:numerics}) use exact diagonalization of the minimal
model, open transverse-field Ising chains~\cite{Pfeuty1970} of up to 12 sites, and harmonic
chains of up to four modes.

% ---------------------------------------------------------------------------------------
\section{Setting and free operations}
\label{sec:setting}

\subsection{Energy booking}

Alice holds a finite-dimensional system $A$ and Bob a system $B$, of dimensions $d_A$ and
$d_B$; the Hamiltonian $H$ acts on $A\otimes B$. Following Hotta~\cite[Sec.~3]{Hotta2011review},
we book the interaction on Bob's side. The part of $H$ supported on $A$ alone,
\begin{equation}
H_A := \tfrac{1}{d_B}\Tr_B(H)\otimes I_B ,
\label{eq:HA}
\end{equation}
is Alice's energy, and $\HBl:=H-H_A$ is Bob's local energy operator; it contains the
interaction. An operator-Schmidt decomposition gives
\begin{equation}
\HBl = I_A\otimes B_0 + \sum_{j=1}^{m} X_j\otimes Y_j ,
\label{eq:schmidt}
\end{equation}
with the $X_j$ traceless, Hermitian, and linearly independent, and the $Y_j$ Hermitian and
linearly independent. We call the $X_j$ Alice's \emph{bond operators}.

\subsection{Free operations}

\begin{itemize}
\item[(O1)] Alice applies an instrument $\{K_\mu\}$ on $A$, with
$\sum_\mu K_\mu^\dagger K_\mu=I$, whose average channel
$\Abar(\cdot)=\sum_\mu K_\mu(\cdot)K_\mu^\dagger$ satisfies
\begin{equation}
(\Abar^\dagger\otimes\id)(\HBl)=\HBl
\;\Longleftrightarrow\;
\Abar^\dagger(X_j)=X_j \;\;\forall j ,
\label{eq:ens}
\end{equation}
where $\Abar^\dagger(Z)=\sum_\mu K_\mu^\dagger Z K_\mu$. (The equivalence follows from
applying $\Abar^\dagger\otimes\id$ to Eq.~\eqref{eq:schmidt}, using
$\Abar^\dagger(I)=I$ and the linear independence of the $Y_j$.)
\item[(O2)] Alice sends classical messages to Bob.
\item[(O3)] Bob applies any channel $\E$ on $B$, which may depend on the messages he has
received. The energy his device receives is
\begin{equation}
w_B(\E,\rho):=\Tr[\HBl\rho]-\Tr[\HBl(\id\otimes\E)(\rho)] .
\label{eq:wB}
\end{equation}
It can be negative, in which case the device supplies energy.
\item[(O4)] Classical randomness may be mixed in, and classical flags discarded or
relabelled.
\end{itemize}

Condition~\eqref{eq:ens} says that Alice's operation, averaged over its outcomes, leaves
$\langle \HBl\rangle$ unchanged in every state: before Bob learns the outcome, Alice cannot
change his energy. We call such instruments \emph{energy-non-signalling}. A sufficient
condition is that every Kraus operator commute with every bond operator,
$[K_\mu,X_j]=0$; we call these instruments \emph{commuting-Kraus}. Hotta's measurement of
$\sigma_x^A$ in the minimal model is of this kind~\cite[Sec.~3]{Hotta2011review}.

The class is closed under composition: if $\{K_\mu\}$ and, for each $\mu$,
$\{K'_{\nu|\mu}\}$ satisfy Eq.~\eqref{eq:ens}, then
$\sum_{\mu\nu}K_\mu^\dagger K'^{\dagger}_{\nu|\mu}X_jK'_{\nu|\mu}K_\mu=\sum_\mu
K_\mu^\dagger X_j K_\mu=X_j$. It is a subclass of one-way LOCC. Condition~\eqref{eq:ens}
places no limit on the energy Alice's apparatus injects into $H_A$; in the minimal model
that energy is Hotta's $E_A$. The theory tracks Bob's side. Bob's device plays the part of
the battery in resource theories of work extraction: every channel on $B$ is allowed, and
the energy it exchanges with the device is booked.

\subsection{Models}
\label{sec:models}

The \emph{minimal model}~\cite[Sec.~3]{Hotta2011review} has two qubits and, up to an
additive constant,
\begin{equation}
H=h\sigma_z^A+h\sigma_z^B+2k\,\sigma_x^A\sigma_x^B ,\qquad h,k>0 .
\label{eq:minimal}
\end{equation}
Here $\HBl=h\sigma_z^B+2k\sigma_x^A\sigma_x^B$, the single bond operator is
$X=\sigma_x^A$, and $Y=2k\sigma_x^B$. The \emph{transverse-field Ising chain} is the open
chain of $L$ qubits, sites $0,\dots,L-1$,
\begin{equation}
H=-J\sum_{i=0}^{L-2}\sigma_x^i\sigma_x^{i+1}-g\sum_{i=0}^{L-1}\sigma_z^i ,
\label{eq:tfim}
\end{equation}
with $J=1$ and $g\in\{0.5,1,2\}$; the critical point is $g=J$. Alice's bond operators are
the $\sigma_x$ of her sites adjacent to Bob's. The \emph{star model} has $n_A$ Alice qubits
and one Bob qubit,
\begin{equation}
H=h\sum_{i=1}^{n_A+1}\sigma_z^i+k\Big(\sum_{i\in A}\sigma_x^i\Big)\otimes\sigma_x^B ,
\label{eq:star}
\end{equation}
with the single bond operator $X=\sum_{i\in A}\sigma_x^i$. It is the one model here that
violates the hypothesis of Theorem~\ref{thm:freeclass}.

% ---------------------------------------------------------------------------------------
\section{The QET potential}
\label{sec:potential}

For an instrument $\{K_\mu\}$ on $A$ write $\tilde\rho_\mu:=(K_\mu\otimes
I)\rho(K_\mu\otimes I)^\dagger$ for the unnormalized branches. The \emph{QET potential} is
\begin{equation}
\begin{split}
\Wto(\rho):=\sup_{\{K_\mu\},\{\E_\mu\}}\sum_\mu\Big\{&\Tr[\HBl\tilde\rho_\mu]\\
&-\Tr\big[\HBl(\id\otimes\E_\mu)(\tilde\rho_\mu)\big]\Big\} ,
\end{split}
\label{eq:W}
\end{equation}
the supremum running over energy-non-signalling instruments and channels $\E_\mu$ on $B$.
By Eq.~\eqref{eq:ens}, $\sum_\mu\Tr[\HBl\tilde\rho_\mu]=\Tr[\HBl\rho]$, so $\Wto(\rho)$ is
$\Tr[\HBl\rho]$ minus the lowest mean energy Bob can reach by free operations. The trivial
protocol shows $\Wto\ge0$. Restricting the supremum to the trivial instrument gives Bob's
local extractable energy $\Wloc(\rho)$; a strongly locally passive state has $\Wloc=0$
\cite{Frey2014,Alhambra2019}. Restricting it to commuting-Kraus instruments gives $\WCK\le\Wto$.
We call states with $\Wto=0$ \emph{free}: no free protocol delivers energy to Bob from them.
In every model here the ground state is not free. The vacuum is the resource.

\subsection{A closed form for commuting bonds}

Suppose the bond operators commute, with joint spectral decomposition
$X_j=\sum_a x_j(a)\Pi_a$, where the $\Pi_a$ are orthogonal projectors summing to $I_A$ and
the joint eigenvalue vectors $x(a)\in\mathbb{R}^m$ are distinct. Equation~\eqref{eq:schmidt}
becomes
\begin{equation}
\HBl=\sum_a\Pi_a\otimes H^{(a)} ,\qquad H^{(a)}:=B_0+\sum_j x_j(a)Y_j .
\label{eq:Heff}
\end{equation}
$H^{(a)}$ is the Hamiltonian Bob faces when Alice's bonds sit in block $a$: the local
effective Hamiltonian of Ref.~\cite{Xie2025}, which was introduced there for a product state.

\begin{proposition}[commuting-Kraus potential]
\label{prop:closed}
For commuting bonds,
\begin{equation}
\WCK(\rho)=\sum_a\Big\{\Tr[H^{(a)}\sigma_a]-p_a\lmin(H^{(a)})\Big\},
\label{eq:closed}
\end{equation}
where $\sigma_a:=\Tr_A[(\Pi_a\otimes I)\rho]$ and $p_a:=\Tr\sigma_a$. The value is attained
by Alice's projective measurement $\{\Pi_a\}$ followed by Bob's reset of $B$ to a ground
state of $H^{(a)}$.
\end{proposition}

\begin{proof}
A commuting-Kraus $K_\mu$ commutes with every $\Pi_a$, since each $\Pi_a$ is a polynomial in
the $X_j$. By Eq.~\eqref{eq:Heff} and because $\E_\mu$ preserves trace and positivity, Bob's
final energy in branch $\mu$ is
$\sum_a\Tr[H^{(a)}\E_\mu(\Tr_A[(\Pi_a\otimes I)\tilde\rho_\mu])]\ge
\sum_a\lmin(H^{(a)})\Tr[(K_\mu^\dagger\Pi_aK_\mu\otimes I)\rho]$. Summing over $\mu$ and using
$\sum_\mu K_\mu^\dagger\Pi_aK_\mu=\sum_\mu K_\mu^\dagger K_\mu\Pi_a=\Pi_a$ bounds the total
final energy below by $\sum_ap_a\lmin(H^{(a)})$, while the initial energy is
$\sum_a\Tr[H^{(a)}\sigma_a]$. The stated protocol leaves branch $a$ with energy
$p_a\lmin(H^{(a)})$, so it attains the bound.
\end{proof}

In the minimal model the blocks are the eigenspaces $a=\pm$ of $\sigma_x^A$,
$H^{(\pm)}=h\sigma_z^B\pm2k\sigma_x^B$, and $\lmin(H^{(\pm)})=-\sqrt{h^2+4k^2}$ for both.
In the ground state~\cite[Sec.~3]{Hotta2011review}, $\langle\sigma_z^B\rangle=-h/\sqrt{h^2+k^2}$
and $\langle\sigma_x^A\sigma_x^B\rangle=-k/\sqrt{h^2+k^2}$, so
$\bra{g}\HBl\ket{g}=-(h^2+2k^2)/\sqrt{h^2+k^2}$. Hence
\begin{equation}
\WCK(\ket{g})=\sqrt{h^2+4k^2}-\frac{h^2+2k^2}{\sqrt{h^2+k^2}} .
\label{eq:EBmax}
\end{equation}
Since $(h^2+4k^2)(h^2+k^2)=(h^2+2k^2)^2+h^2k^2$, this equals Hotta's
$E_B^{\max}=\frac{h^2+2k^2}{\sqrt{h^2+k^2}}\big[\sqrt{1+h^2k^2/(h^2+2k^2)^2}-1\big]$
\cite[Eq.~(11)]{Hotta2011review}. Alice holds a qubit, so commuting-Kraus instruments are her
whole free class (Theorem~\ref{thm:freeclass}), and $\Wto(\ket g)=E_B^{\max}$. Hotta's
protocol is therefore optimal among all free protocols, and $E_B^{\max}$ is the resource it
spends. Hotta's injected energy is $E_A=h^2/\sqrt{h^2+k^2}$ \cite[Eq.~(8)]{Hotta2011review},
so with $r:=k/h$ and $x:=1+2r^2$,
\begin{equation}
\frac{E_B^{\max}}{E_A}=\sqrt{x^2+r^2}-x<\frac{r^2}{2x}<\frac14 ,
\label{eq:quarter}
\end{equation}
using $\sqrt{x^2+r^2}<x+r^2/(2x)$. In the minimal model Bob never receives a quarter of what
Alice injects; the ratio approaches $1/4$ only as $k/h\to\infty$ (Fig.~\ref{fig:minimal}).
Equation~\eqref{eq:quarter} is an elementary consequence of Hotta's two formulas.

\subsection{Monotonicity}

\begin{theorem}[ledger monotonicity]
\label{thm:ledger}
For every state $\rho$:
\begin{itemize}
\item[(a)] for every energy-non-signalling instrument on $A$, with branches $\rho_\mu$ of
probability $p_\mu$,
\begin{equation}
\sum_\mu p_\mu\Wto(\rho_\mu)\le\Wto(\rho);
\label{eq:mono-a}
\end{equation}
\item[(b)] for every channel $\E$ on $B$,
\begin{equation}
\Wto\big((\id\otimes\E)(\rho)\big)+w_B(\E,\rho)\le\Wto(\rho).
\label{eq:mono-b}
\end{equation}
\end{itemize}
Both statements hold for $\WCK$ when (a) is restricted to commuting-Kraus instruments.
\end{theorem}

The proof is in Appendix~\ref{app:proofs}. Inequality~\eqref{eq:mono-a} is strong
monotonicity, i.e., monotonicity on average over branches~\cite{ChitambarGour2019}. Because
$\Wto$ is a supremum of linear functionals, it is convex, so it is also non-increasing under
the average channel and under mixing.

\begin{corollary}[device ledger]
\label{cor:ledger}
Let $D$ be the energy delivered to Bob's device so far. Along any free protocol, the
expected value of $\Wto(\text{state})+D$ does not increase. Hence the energy any free
protocol delivers to Bob's device is at most $\Wto(\rho)$; $\Wto$ does not increase under
free operations that draw no energy from the device; and a free operation can raise $\Wto$
only by drawing energy from the device, and by at most the amount drawn.
\end{corollary}

In the minimal model the ledger reads as follows. Alice's measurement leaves
$\Wto=E_B^{\max}$ on average (Theorem~\ref{thm:linear} below makes this an equality). After
the measurement the energy is locally extractable on each branch, because Alice's state
lies in a single block. Bob's rotation delivers $E_B^{\max}$ and leaves branches with
$\Wto=0$. Corollary~\ref{cor:ledger} is not a statement about $\Wto$ alone. Bob may excite
his own system; that raises $\Wto$, but only by the energy he put in.

% ---------------------------------------------------------------------------------------
\section{The free class}
\label{sec:freeclass}

\begin{theorem}[free class]
\label{thm:freeclass}
Let the bond operators commute, with blocks $\Pi_a$ and joint eigenvalue vectors $x(a)$.
Every energy-non-signalling instrument on $A$ is commuting-Kraus if and only if every $x(a)$
is an extreme point of $\mathrm{conv}\{x(a')\}_{a'}$, i.e., the $x(a)$ are in convex
position.
\end{theorem}

The key identity is the following. Set
$M_{ab}:=\sum_\mu\Pi_bK_\mu^\dagger\Pi_aK_\mu\Pi_b\ge0$. Sandwiching completeness and
Eq.~\eqref{eq:ens} between $\Pi_b$'s gives
\begin{equation}
\sum_aM_{ab}=\Pi_b ,\qquad \sum_ax(a)M_{ab}=x(b)\,\Pi_b .
\label{eq:blocks}
\end{equation}
For a unit vector $v$ in the range of $\Pi_b$, the numbers $p_a=\bra vM_{ab}\ket v$ thus form
a probability distribution with barycentre $x(b)$. If $x(b)$ is extreme, $p_a=0$ for
$a\neq b$. Hence $M_{ab}=0$, i.e., $\Pi_aK_\mu\Pi_b=0$, and every $K_\mu$ is block diagonal.
Conversely, suppose $x(b)=\sum_{a\neq b}\lambda_ax(a)$ for a probability vector $\lambda$.
Choose unit vectors $\ket b$ in the range of $\Pi_b$ and $\ket a$ in the range of each
$\Pi_a$. Then
\begin{equation}
K_0=I-\proj b ,\qquad K_a=\sqrt{\lambda_a}\,\ket a\!\bra b
\label{eq:witness}
\end{equation}
is complete and satisfies Eq.~\eqref{eq:ens}, since
$\sum K^\dagger X_jK=X_j-x_j(b)\proj b+\sum_a\lambda_ax_j(a)\proj b=X_j$. Yet
$\Pi_aK_a\Pi_b\neq0$. The witness~\eqref{eq:witness} is non-unital:
$\Abar(I)=I-\proj b+\sum_a\lambda_a\proj a\neq I$. That is unavoidable. For a unital
instrument, $\Abar^\dagger$ is a unital, trace-preserving channel, and its fixed points form
the commutant of the Kraus operators~\cite{AGG2002}. So a unital energy-non-signalling
instrument is commuting-Kraus whatever the spectrum, and convex position is exactly what the
non-unital instruments need. The complete proof is in Appendix~\ref{app:proofs}.

Which models are in convex position? A single-qubit $A$ always is. Commuting traceless
Hermitian $2\times2$ matrices are proportional to one $\bm n\cdot\bm\sigma$, so the points
are $\pm v$. Every region of the Ising chain~\eqref{eq:tfim} is too, since its bond operators
are $\sigma_x$'s whose joint eigenvalues are the vertices of a box. The star
model~\eqref{eq:star} with $n_A\ge2$ is not: $X=\sum_i\sigma_x^i$ has the eigenvalues
$n_A,n_A-2,\dots,-n_A$, and all but the two outermost are interior.

\subsection{A certificate}

Everything in Eq.~\eqref{eq:blocks} depends on the instrument only through its average
channel, $M_{ab}=\Pi_b\Abar^\dagger(\Pi_a)\Pi_b$, and Eq.~\eqref{eq:ens} is linear in the
Choi operator $J$ of $\Abar$. Maximizing the weight outside the block-diagonal,
\begin{equation}
w_{\mathrm{off}}(J):=\sum_{a\neq b}\Tr M_{ab},
\label{eq:sdp}
\end{equation}
subject to $J\ge0$, trace preservation, and Eq.~\eqref{eq:ens}, is therefore a semidefinite
program. Its optimum is zero exactly when every energy-non-signalling instrument is
commuting-Kraus. The optimum is in fact
$\sum_{b\,\text{non-extreme}}\operatorname{rank}\Pi_b$. It is at most that, because by
Eq.~\eqref{eq:blocks} the extreme blocks contribute nothing and block $b$ contributes at most
$\Tr\Pi_b$. Applying the witness~\eqref{eq:witness} to every vector of an orthonormal basis of
every non-extreme block attains it; take $\lambda_b=0$, which is possible because a
non-extreme $x(b)$ lies in the hull of the other points. Table~\ref{tab:freeclass} lists
the solver's optima. They agree with the analytic verdict on every row and equal the total
rank of the non-extreme blocks: 2 and 6 on the two star models.

\begin{table}[t]
\caption{The free class of each model. Ranks are those of the joint eigenprojectors
$\Pi_a$ of Alice's bond operators. $w_{\rm off}^{\rm SDP}$ is the solver optimum of
Eq.~\eqref{eq:sdp}. The last column is the total rank of the non-extreme blocks, which
equals the exact optimum. Minimal-model rows use $(h,k)=(1,\tfrac12)$ with Alice at qubit 0
and $(0.3,2)$ with Alice at qubit 1; Alice's sites are listed. Ising rows use $g=1$. Star rows use $(h,k)=(1,\tfrac12)$.}
\label{tab:freeclass}
\begin{ruledtabular}
\begin{tabular}{lllcrr}
Model & Alice & Ranks & Convex & $w_{\rm off}^{\rm SDP}$ & Exact \\
\hline
Minimal & $\{0\}$ & 1,1 & yes & $2\times10^{-16}$ & 0\\
Minimal & $\{1\}$ & 1,1 & yes & $1\times10^{-15}$ & 0\\
Ising $L=5$ & $\{0\}$ & 1,1 & yes & $1\times10^{-15}$ & 0\\
Ising $L=5$ & $\{1,2\}$ & 1,1,1,1 & yes & $-8\times10^{-11}$ & 0\\
Ising $L=6$ & $\{2,3\}$ & 1,1,1,1 & yes & $-8\times10^{-11}$ & 0\\
Star $n_A=2$ & $\{0,1\}$ & 1,2,1 & no & 2.000 & 2\\
Star $n_A=3$ & $\{0,1,2\}$ & 1,3,3,1 & no & 6.000 & 6\\
\end{tabular}
\end{ruledtabular}
\end{table}

\subsection{Off convex position}
\label{sec:offconvex}

Take the star model with $n_A=2$, $(h,k)=(1,\tfrac12)$, and its ground state. There
$\WCK(\ket g)=0.036634$. Apply the witness~\eqref{eq:witness}, whose energy-signalling
defect is $9\times10^{-16}$, then the optimal commuting-Kraus protocol on each branch. Bob
receives $0.225331$, which is 6.15 times as much. With only the witness outcome, Bob's local
extraction is already $0.188697$ (5.15 times $\WCK$). Hence $\Wto(\ket g)\ge0.225331>\WCK(\ket g)$
there. The closed form~\eqref{eq:closed} is not the potential, and it is not a monotone under
the free class: its average over the witness's branches exceeds its initial value by
$0.188697$. For $n_A=3$ a witness exists, but this particular one extracts nothing measurable
($6\times10^{-14}$). The optimal free instrument off convex position is not known.

% ---------------------------------------------------------------------------------------
\section{Linearity and two-way communication}
\label{sec:twoway}

\begin{theorem}[linearity and the exact ledger]
\label{thm:linear}
For commuting bonds, $\WCK(\rho)=\Tr[W_{\rm op}\rho]$ with
\begin{equation}
W_{\rm op}:=\sum_a\Pi_a\otimes\big(H^{(a)}-\lmin(H^{(a)})I\big)=\HBl+G ,
\label{eq:Wop}
\end{equation}
where $G:=-\sum_a\lmin(H^{(a)})\,\Pi_a\otimes I_B$ is supported on $A$. Consequently:
\begin{itemize}
\item[(a)] for every instrument $\{\mathcal L_\nu\}$ on $B$, with branches $\rho_\nu$ of
probability $p_\nu$, $\sum_\nu p_\nu\WCK(\rho_\nu)=\WCK(\rho)-w_B$, where
$w_B=\Tr[\HBl\rho]-\sum_\nu\Tr[\HBl(\id\otimes\mathcal L_\nu)(\rho)]$ is the energy the
instrument delivers to Bob's device;
\item[(b)] for every commuting-Kraus instrument on $A$,
$\sum_\mu p_\mu\WCK(\rho_\mu)=\WCK(\rho)$.
\end{itemize}
If the $x(a)$ are in convex position, then $\Wto=\WCK$ and (a) and (b) hold for all free
operations: the ledger of Corollary~\ref{cor:ledger} is conserved exactly.
\end{theorem}

\begin{proof}
Linearity follows from Eqs.~\eqref{eq:closed} and~\eqref{eq:Heff}. For (a), the sum over
branches of $\Tr[G(\id\otimes\mathcal L_\nu)(\rho)]$ equals $\Tr[G\rho]$ because $G$ acts on
$A$ alone and $\sum_\nu\mathcal L_\nu$ is trace preserving; the $\HBl$ part changes by
$-w_B$. For (b), commuting-Kraus instruments satisfy
$\sum_\mu K_\mu^\dagger\Pi_aK_\mu=\Pi_a$, which preserves $G$, and Eq.~\eqref{eq:ens},
which preserves $\HBl$. The last sentence follows from Theorem~\ref{thm:freeclass}.
\end{proof}

\begin{corollary}[two-way communication gives Bob nothing]
\label{cor:twoway}
Consider protocols of finitely many rounds. In each round either Alice applies a
commuting-Kraus instrument on $A$ or Bob applies an instrument on $B$; the outcome is
broadcast, and every choice may depend on all earlier outcomes. The expected energy
delivered to Bob's device is $\WCK(\rho)-\mathbb{E}[\WCK(\rho_{\rm fin})]\le\WCK(\rho)$,
and one-way communication attains it. In convex position this holds for all free
instruments on $A$, and the two-way value equals $\Wto(\rho)$.
\end{corollary}

\begin{proof}
By Theorem~\ref{thm:linear}, the expectation of $\WCK(\text{state})+D$ is the same after
every round. It starts at $\WCK(\rho)$ and ends at no less than $D$, since $\WCK\ge0$.
Proposition~\ref{prop:closed} attains it in one round.
\end{proof}

The mechanism is linearity. Because $\WCK$ is linear, what Bob learns by measuring cannot
raise its expectation, and Alice's commuting-Kraus instruments conserve it whatever they are
conditioned on. Bob's own instruments pay for every rise with an equal injection.

The scope matters. Off convex position, part (b) fails for free instruments outside the
commuting class (Sec.~\ref{sec:offconvex}), and the failure already occurs with one-way
communication. Corollary~\ref{cor:twoway} then says nothing about the true potential
$\Wto$. Whether two-way communication can beat one-way communication when the bond spectrum
is not in convex position is open. We also restrict to finitely many rounds; the round
number of LOCC protocols is itself a resource~\cite{ChitambarGour2019}.

\begin{remark}[crediting both sides]
Two-way protocols can deliver more energy in total if a device on Alice's side is also
credited. To see this, book the interaction half to each side, which leaves
Eq.~\eqref{eq:ens} and hence Alice's free class unchanged. In each round one party measures
its finest free instrument and the other cashes out optimally, with the measuring party
chosen to maximize the total. The total delivered
then grows with the number of rounds $r$ (Table~\ref{tab:twoway}), reaching $20.5\,\Wto$ at
$r=3$ for $(h,k)=(1,\tfrac12)$. But the measurements inject exactly as much more. The net
energy taken from the system, delivered minus injected, is $-0.8219$ for $r=1,2,3$. It is
round-independent to $2\times10^{-13}$ in all seven cases, and negative, as it must be for a
ground state, which is passive~\cite{PuszWoronowicz1978}. The extra delivery is the measuring
apparatus's energy passing through the system.
\end{remark}

\begin{table}[t]
\caption{Two-way protocols with both sides' devices credited. $\Wto$ is the one-way
potential; the next three columns give the total delivered after $r=1,2,3$ rounds; the net
is delivered minus injected, which is the same for every $r$. Minimal-model rows are labelled
by $(h,k)$; Ising rows have $g=1$ and Alice at site 0. The $L=7,8$ rows were run to two
rounds.}
\label{tab:twoway}
\begin{ruledtabular}
\begin{tabular}{lccccc}
Case & $\Wto$ & $r=1$ & $r=2$ & $r=3$ & Net\\
\hline
$(1,\tfrac12)$ & 0.0726 & 0.0726 & 0.7797 & 1.4868 & $-0.8219$\\
$(0.3,2)$ & 0.0110 & 0.0110 & 0.0334 & 0.0559 & $-0.0335$\\
$(2,2)$ & 0.2295 & 0.2295 & 1.1239 & 2.0183 & $-1.1847$\\
$L=5$ & 0.1447 & 0.1447 & 0.7410 & 1.3832 & $-0.7129$\\
$L=6$ & 0.1546 & 0.1546 & 0.7476 & 1.3861 & $-0.7005$\\
$L=7$ & 0.1620 & 0.1620 & 0.7530 & --- & $-0.6915$\\
$L=8$ & 0.1677 & 0.1677 & 0.7575 & --- & $-0.6848$\\
\end{tabular}
\end{ruledtabular}
\end{table}

% ---------------------------------------------------------------------------------------
\section{The surplus and the negativity}
\label{sec:surplus}

A natural candidate for ``what classical communication buys'' is the surplus
$Q:=\Wto-\Wloc\ge0$. It is not a monotone.

\begin{proposition}[the surplus]
\label{prop:surplus}
Let the bonds commute, let Bob hold all of $B$, and restrict Alice to commuting-Kraus
instruments. Then:
\begin{itemize}
\item[(a)] $Q$ does not increase on average under Alice's instruments, and does not
decrease under channels on $B$.
\item[(b)] $Q$ can increase under a local unitary on $A$, an LOCC operation outside the free
class. In the minimal model, $\rho_1=\proj{+}\otimes\proj0$ has $Q=0$. The unitary
$U_A:\ket+\mapsto\ket0$ maps it to $\rho_2=\proj0\otimes\proj0$, which has
$Q=\sqrt{h^2+4k^2}-h$.
\item[(c)] $Q$ can increase under a free channel on $B$ that delivers energy to Bob's device.
The state $\rho_3=\tfrac12(\proj+\otimes\proj0+\proj-\otimes\proj1)$ has $Q=0$. Bob's reset
of $B$ to $\ket1$ delivers $h$ and yields $\rho_5=\tfrac{I}{2}\otimes\proj1$, which has
$Q=\sqrt{h^2+4k^2}-h$.
\end{itemize}
Here $\ket\pm$ are the eigenvectors of $\sigma_x$, and $\sigma_z\ket0=\ket0$,
$\sigma_z\ket1=-\ket1$.
\end{proposition}

The proof is in Appendix~\ref{app:proofs}. $Q$ is therefore neither non-increasing nor
non-decreasing under the free operations, and it is not an LOCC monotone. In (c) Bob
extracts energy and $Q$ rises, so $Q$ is not what Bob spends; $\Wto$ is. In (c) $\Wto$ falls
from $\sqrt{h^2+4k^2}$ to $\sqrt{h^2+4k^2}-h$, by exactly the delivered $h$, as
Theorem~\ref{thm:linear}(a) requires. At $(h,k)=(1,\tfrac12)$ the rises in (b) and (c) are
both $\sqrt2-1=0.41421$, and the code reproduces them. Over 60 random commuting-Kraus
instruments, $Q$ never rose by more than $6\times10^{-14}$. Over 60 random channels on $B$ it
never fell by more than $9\times10^{-14}$.

The negativity $N(\rho)=(\|\rho^{T_A}\|_1-1)/2$ is an LOCC monotone~\cite{VidalWerner2002}, and
hence a monotone under the free operations, which are a subclass of LOCC. It is not what Bob
spends, for three reasons. In the minimal model Alice's measurement leaves product branches,
with $N=0$, that still carry $\Wto=E_B^{\max}$. The state $\rho_2$ above is a product state
with $Q>0$. On the chains, the two-site negativity $N(\rho_{AB})$ vanishes for $d\ge3$ while
Bob's yield stays positive (Fig.~\ref{fig:chain}); the correlations Bob uses at those
distances are carried by the rest of the chain. Hotta derived relations between teleported
energy and ground-state entanglement for the minimal model~\cite{Hotta2010}.

% ---------------------------------------------------------------------------------------
\section{A single-qubit receiver}
\label{sec:single}

On the chains Bob holds one site $b$ at distance $d$ from Alice. An operation on $b$
changes only the terms of $H$ that touch $b$. We evaluate energies with an operator $H_R$ on
$R=\{b-1,b,b+1\}$ that contains every such term: the part of the three sites' local energy
densities that is supported on $R$. Let $\tau_\mu$
be the conditional state of $R$ after Alice's outcome $\mu$. Bob's yield from unitaries
$U_\mu$ on $b$ is $E_B=\sum_\mu
p_\mu\big(\Tr[H_R\tau_\mu]-\Tr[H_RU_\mu\tau_\mu U_\mu^\dagger]\big)$.

\subsection{The exact unitary yield}

Write $H_R=I_b\otimes C_0+\sum_{i=1}^3\sigma_i\otimes C_i$ and let
$T_{ij}:=\Tr[(\sigma_j\otimes C_i)\tau]$, a real $3\times3$ matrix.

\begin{proposition}[Salvia, De Palma, and Giovannetti~\cite{Salvia2023}]
\label{prop:salvia}
The largest energy decrease of $\tau$ under unitaries on the qubit $b$ is
\begin{equation}
E^{\rm unit}(\tau)=\Tr T+s_1+s_2-\sgn(\det T)\,s_3 ,
\label{eq:salvia}
\end{equation}
where $s_1\ge s_2\ge s_3$ are the singular values of $T$.
\end{proposition}

This is the closed formula for the local ergotropy of a two-level system in Sec.~III~A
of Ref.~\cite{Salvia2023}, their Eq.~(19). Their matrix $M$ (their Eq.~(18)) is $-T^{\top}$ in
the present notation. The argument is short. $U^\dagger\sigma_iU=\sum_jR_{ij}\sigma_j$ with
$R\in SO(3)$, and every $R\in SO(3)$ arises. So the energy after the unitary is
$\Tr[C_0\Tr_b\tau]+\sum_{ij}R_{ij}T_{ij}$. Minimizing over $SO(3)$ is the orthogonal
Procrustes problem with a determinant constraint~\cite{Kabsch1976}.

We apply Eq.~\eqref{eq:salvia} branch by branch,
$E_B^{\rm unit}=\sum_\mu p_\mu E^{\rm unit}(\tau_\mu)$. We used 96 chain configurations:
$L\in\{6,8,10,12\}$, $g\in\{0.5,1,2\}$, Alice at site 0, and Bob at every distance. On all of
them $E_B^{\rm unit}$ equals the yield of the chain protocol's optimized conditional
rotation~\cite{Hotta2009} to $4.4\times10^{-16}$ in absolute energy. The one-parameter
rotation family of the protocol is therefore optimal among all unitaries on $b$. On the
minimal model Eq.~\eqref{eq:salvia} reproduces $E_B^{\max}$ to $10^{-12}$.

Unitaries are not Bob's best free operations in the ordered phase. At $g=0.5$ and $d=1$, a
channel on $b$ beats every unitary. At $L=6$ the Choi-matrix semidefinite program of
Ref.~\cite{Alhambra2019} certifies $9.697\times10^{-3}$ against $1.654\times10^{-3}$, a factor
5.86. A numerically optimized channel, a lower bound on the optimum, gives factors of 3.2,
5.86, 8.5, 10.5, 11.7, and 12.4 for $L=5,\dots,10$. At $g=1$ and $g=2$ the rotation attains
the single-site channel optimum to $2\times10^{-7}$, relative, on configurations with
$E_B>10^{-6}$. That local channels can beat local unitaries is known~\cite{Alhambra2019,Salvia2023};
the Ising values are measurements, nothing more.

\subsection{Entropic bounds, and why they are loose}

\begin{proposition}[entropic bounds]
\label{prop:entropic}
Let $\varepsilon_1\le\varepsilon_2\le\dots$ be the eigenvalues of $H_R$ and
$\Hmin(\tau):=-\ln\lambda_{\max}(\tau)$.
\begin{itemize}
\item[(a)] For unitaries on $R$ or on any part of it,
\begin{equation}
E_B\le\sum_\mu p_\mu\Big(\Tr[H_R\tau_\mu]-\kappa\big(H_R,\Hmin(\tau_\mu)\big)\Big),
\label{eq:kyfan}
\end{equation}
where $\kappa(H,s):=\min\{\Tr[H\omega]:\omega\ge0,\Tr\omega=1,\omega\le e^{-s}I\}$. This
minimum equals $e^{-s}\sum_{i\le n}\varepsilon_i+(1-ne^{-s})\varepsilon_{n+1}$, with
$n=\lfloor e^{s}\rfloor$.
\item[(b)] Suppose the average $\rho_R=\sum_\mu p_\mu\tau_\mu$ is passive for Bob's unitaries on
$b$, and Alice's outcome is binary with $p_\pm=\tfrac12$. Then
\begin{equation}
E_B\le2c\Big(e^{-\Hmin(X|R)}-\tfrac12\Big),
\label{eq:guess}
\end{equation}
where $e^{-\Hmin(X|R)}$ is the probability of guessing Alice's outcome $X$ from
$R$~\cite{KRS2009}, and $c:=\tfrac12\max_U\spread\big(H_R-(U^\dagger\otimes
I)H_R(U\otimes I)\big)$ over unitaries $U$ on $b$.
\end{itemize}
\end{proposition}

Bound~(a) holds because $U\tau U^\dagger$ has the spectrum of $\tau$ and so is feasible for
$\kappa$. Bound~(b) combines passivity with the H\"older inequality and the Helstrom formula
(Appendix~\ref{app:proofs}). The passivity hypothesis holds when the global state is strongly
locally passive and Alice's instrument is energy-non-signalling; it was checked on every
configuration.

Bound~(a) is attained when the conditional states are pure and Bob acts on all of $R$. Then
$\Hmin=0$ and $\kappa=\varepsilon_1$, so it is the pure-state ergotropy. This happens in the
minimal model,
where the ratio is 1 to $1.5\times10^{-11}$. Saturation there reflects purity, not a tight
entropic principle. On the chains, where $\tau_\mu$ is mixed, the bound is loose: $E_B$ is at
most $0.041$ of it. Replacing $\Hmin$ by its smoothed version~\cite{Tomamichel2016} did not
tighten it on any of 126 chain configurations. Bound~(b) says that Bob extracts only what his
region knows about Alice's outcome, and it vanishes exactly when the outcome is unguessable.
But it is loose too: on the 96 configurations, $E_B$ is at most $0.0080$ of it.

The looseness has an identifiable cause. Alice's outcome moves $\tau_\mu$ by an $O(1)$ trace
distance, but in a nearly energy-neutral direction. The tight value~\eqref{eq:salvia} comes
from a cancellation, inside the nuclear norm of $T$, between Bob's field row and the
correlation rows, and norm-based bounds discard it. An entropic quantity that retains the
cancellation is not known to us. Itoh, Masaki, and Matsueda derived information-thermodynamic
bounds of a related kind for measurement-and-feedback extraction from pure
states~\cite{Itoh2026}.

\subsection{A daemonic gain}

Francica \emph{et al.}~\cite{Francica2017} define the daemonic ergotropy. It is the energy
extractable from a system by unitaries conditioned on the outcome of a measurement on a
correlated ancilla (their Eq.~(3)); the daemonic gain is its excess over the unconditioned
ergotropy (their Eq.~(5)). Bernards \emph{et al.} extended the notion to generalized
measurements~\cite{Bernards2019}. Read $R$ as the system and the rest of the chain as the
ancilla, and restrict the daemon's measurements to free instruments. Bob's optimal unitary
yield is then a daemonic ergotropy. Without the message, Bob's unitary yield vanishes by
strong local passivity, so $E_B$ is entirely daemonic gain. Itoh \emph{et al.} make the same
connection between feedback extraction and daemonic ergotropy~\cite{Itoh2026}. The resource
theory adds a constraint on the daemon, not a new bound.

% ---------------------------------------------------------------------------------------
\section{Numerical checks}
\label{sec:numerics}

\begin{figure}[t]
\centering
\begin{tikzpicture}
\begin{axis}[
  width=\columnwidth, height=0.64\columnwidth,
  xmode=log, xmin=0.02, xmax=80, ymin=0, ymax=0.27,
  xlabel={$k/h$}, ylabel={$\Wto(\ket g)/E_A$},
  tick label style={font=\footnotesize}, label style={font=\small},
  legend pos=south east, legend cell align=left,
  legend style={font=\footnotesize, draw=none, fill=none},
  ytick={0,0.05,0.1,0.15,0.2,0.25},
  yticklabel style={/pgf/number format/fixed, /pgf/number format/precision=2},
]
\addplot[thin, color=black] table[x=k_over_h, y=rate_closed_form] {data/fig1_closed.dat};
\addlegendentry{Hotta's $E_B^{\max}/E_A$}
\addplot[only marks, mark=*, mark size=0.7pt, color=blue!70!black]
  table[x=k_over_h, y=W_arrow_over_E_A] {data/fig1_grid.dat};
\addlegendentry{Eq.~\eqref{eq:closed}, $33\times33$ grid}
\addplot[dashed, color=gray, domain=0.02:80, samples=2] {0.25};
\addlegendentry{$1/4$}
\end{axis}
\end{tikzpicture}
\caption{The minimal model. The potential~\eqref{eq:closed} of the ground state, computed
from the Hamiltonian on a logarithmic $33\times33$ grid ($h\in[0.05,5]$, $k\in[0.1,4]$) and
divided by Alice's injected energy $E_A$, collapses onto a function of $k/h$. It agrees with
Hotta's closed form to $3.7\times10^{-15}$ and stays below $1/4$ [Eq.~\eqref{eq:quarter}]; the
grid maximum is $0.249978$ at $(h,k)=(0.05,4)$.}
\label{fig:minimal}
\end{figure}

\begin{figure}[t]
\centering
\begin{tikzpicture}
\begin{groupplot}[
  group style={group size=1 by 2, vertical sep=1.05cm},
  width=\columnwidth, height=0.56\columnwidth,
  ymode=log, xmin=0.5, xmax=9.5, ymin=1e-9, ymax=3,
  xtick={1,...,9},
  tick label style={font=\footnotesize}, label style={font=\small},
  title style={font=\small, yshift=-1.2ex},
  unbounded coords=discard, filter discard warning=false,
  ylabel={energy},
]
\nextgroupplot[title={$g=1$ (critical)}, legend to name=chainlegend, legend columns=2,
  legend cell align=left,
  legend style={font=\scriptsize, draw=none, /tikz/every even column/.append style={column sep=8pt}}]
\addplot[only marks, mark=*, mark size=1.6pt, color=blue!70!black]
  table[x=d, y=E_B] {data/fig2_L10_g1.dat};
\addlegendentry{$E_B$ (optimal rotation)}
\addplot[only marks, mark=square, mark size=2.6pt, color=red!70!black]
  table[x=d, y=W_site] {data/fig2_L10_g1.dat};
\addlegendentry{$W_{\rm site}$ (best channel on $b$)}
\addplot[only marks, mark=triangle*, mark size=2pt, color=gray]
  table[x=d, y=bound_region] {data/fig2_L10_g1.dat};
\addlegendentry{bound~\eqref{eq:kyfan}}
\addplot[only marks, mark=diamond*, mark size=2pt, color=green!50!black]
  table[x=d, y=negativity] {data/fig2_L10_g1.dat};
\addlegendentry{$N(\rho_{AB})$}
\addplot[dashed, color=black, domain=0.5:9.5, samples=2] {0.175927};
\addlegendentry{$\Wto$, Bob $=$ sites $1$--$9$}
\nextgroupplot[title={$g=0.5$ (ordered)}, xlabel={distance $d$ between Alice and Bob}]
\addplot[only marks, mark=*, mark size=1.6pt, color=blue!70!black]
  table[x=d, y=E_B] {data/fig2_L10_g0.5.dat};
\addplot[only marks, mark=square, mark size=2.6pt, color=red!70!black]
  table[x=d, y=W_site] {data/fig2_L10_g0.5.dat};
\addplot[only marks, mark=triangle*, mark size=2pt, color=gray]
  table[x=d, y=bound_region] {data/fig2_L10_g0.5.dat};
\addplot[only marks, mark=diamond*, mark size=2pt, color=green!50!black]
  table[x=d, y=negativity] {data/fig2_L10_g0.5.dat};
\addplot[dashed, color=black, domain=0.5:9.5, samples=2] {0.112409};
\end{groupplot}
\node[anchor=north, yshift=-3pt] at (group c1r2.outer south) {\pgfplotslegendfromname{chainlegend}};
\end{tikzpicture}
\caption{An Ising chain of $L=10$ sites; Alice at site 0 measures $\sigma_x$, and Bob holds
the site at distance $d$. $E_B$ is the yield of the optimized conditional rotation, which
equals the exact unitary yield~\eqref{eq:salvia}. $W_{\rm site}$ is the best single-site
channel, found numerically. The triangles show bound~\eqref{eq:kyfan} on the region $R$.
$N(\rho_{AB})$ is the two-site negativity of the ground state; zeros are not drawn. The
dashed line is the potential with Bob holding the whole complement: $0.1759$ at $g=1$ and
$0.1124$ at $g=0.5$, against $E_A=0.851$ and $0.242$. In the ordered phase the best channel
beats the rotation at $d=1$ by a factor of 12.4.}
\label{fig:chain}
\end{figure}

\subsection{The minimal model}

On the grid of Fig.~\ref{fig:minimal}, Eq.~\eqref{eq:closed} is computed from the
Hamiltonian and its ground state. It agrees with Hotta's $E_B^{\max}$ to $3.7\times10^{-15}$,
bound~\eqref{eq:kyfan} is attained to $1.5\times10^{-11}$ relative, and the local
extractable energy is at most $2.3\times10^{-15}$. The last number is strong local passivity,
certified by the semidefinite program of Ref.~\cite{Alhambra2019}.

\subsection{Random free operations}

We applied 12\,500 random free operations to random pure, mixed, ground-state-mixture, and
thermal states. Of these, 10\,550 acted on systems of 2 to 6 qubits, in seven cases: the
minimal model at three values of $(h,k)$, and Ising chains with $L=4$--$6$ and Alice regions
of one or two sites. The operations were commuting-Kraus instruments with 1 to 3 Kraus
operators, channels on $B$, and unitaries on $B$. The other 1\,950 acted on harmonic chains of
2 to 4 modes (Sec.~\ref{sec:gaussian}). There the operations were Gaussian unitaries on $A$
that commute with the bond quadratures, symplectic maps on $B$, and loss on $B$. The worst
violation of Eqs.~\eqref{eq:mono-a} and~\eqref{eq:mono-b}, and of negativity monotonicity,
was $2.1\times10^{-14}$. The energy-signalling defect of the sampled instruments was at most
$1.8\times10^{-14}$.

The checks can fail. In every case a random non-free instrument on $A$ raised $\Wto$, by
0.96 to 7.72 on qubits and by 0.52 to 2.48 on modes. An energy-injecting channel on $B$ raised
it by 2.19 to 7.50, but never by more than the injected energy, to within $9.2\times10^{-15}$.

\subsection{Two-way protocols}

On the minimal model we computed the optimal finite-round two-way value with Bob-only
cash-out. The computation recurses over which party measures in each round, for up to three
rounds, and the result equals $\Wto$ to $2\times10^{-15}$. On chains the recursion would need
Bob's full channel optimum on entangled states, so there the evidence is sampled instead. We
ran 1\,180 random two-way protocols on the minimal model and on Ising chains of 4 to 6 sites.
Each is a sequence of three moves. A move is a commuting-Kraus instrument on $A$ conditioned
on the record so far, an instrument on $B$ whose outcome Alice may use, or a cash-out on $B$.
Across them the ledger of Corollary~\ref{cor:ledger} is conserved to $9.0\times10^{-15}$,
while a non-free instrument on $A$ breaks it by at least $1.08$.

\subsection{The chains}

Figure~\ref{fig:chain} shows an Ising chain of 10 sites; the data cover
$L=4$--$12$ and $g\in\{0.5,1,2\}$. Bob's single-site yield is at most $0.686$ of the
potential with Bob holding the whole complement; the maximum occurs at $g=2$, $d=1$. At
$L=12$ the potential~\eqref{eq:closed} and the region quantities are computed without forming
dense $2^L\times2^L$ matrices, and they agree with the dense computation to $10^{-9}$ where
both run.

\subsection{A Gaussian analogue}
\label{sec:gaussian}

For a harmonic chain $H=\tfrac12(p\cdot p+x\cdot Kx)$ with Bob holding the complement, the
bond terms are $x_aK_{ab}x_b$. Homodyne detection of Alice's bond quadratures $x_a$
commutes with them. Given the outcome vector $\xi$, Bob's conditional state is Gaussian, and
his effective Hamiltonian acquires the linear term $\xi^{\top}K_{AB}x_B$. That lowers its
ground energy by $\tfrac12\xi^{\top}K_{AB}K_{BB}^{-1}K_{AB}^{\top}\xi$. Averaging over $\xi$
gives the analogue of Eq.~\eqref{eq:closed},
\begin{equation}
W=\langle\HBl\rangle-E_0(K_{BB})+\tfrac12\Tr\big[K_{AB}K_{BB}^{-1}K_{AB}^{\top}\Sigma\big],
\label{eq:gauss}
\end{equation}
where $\Sigma$ is the covariance of the bond quadratures and $E_0(K_{BB})$ is the ground
energy of Bob's uncoupled chain. For 2 to 4 modes it agrees with the explicit
conditional-state integral to $6.0\times10^{-16}$, with $W$ ranging from $0.0017$ to $0.045$
over the masses studied. Equation~\eqref{eq:gauss} is the value of the homodyne protocol; we
prove no Gaussian or infinite-dimensional analogue of Theorems~\ref{thm:freeclass}
and~\ref{thm:linear}.

% ---------------------------------------------------------------------------------------
\section{Related work}
\label{sec:related}

\emph{Local extraction.} Strong local passivity was characterized in
Refs.~\cite{Frey2014,Alhambra2019}. Alhambra \emph{et al.} cast local channel extraction as a
semidefinite program. Salvia, De Palma, and Giovannetti introduced the local
ergotropy~\cite{Salvia2023}, whose qubit formula is our Proposition~\ref{prop:salvia}.
Castellano \emph{et al.} extended it by exploiting the free evolution of system and
environment~\cite{Castellano2024}. Bhattacharyya \emph{et al.} studied maps that are not
completely positive~\cite{Bhattacharyya2024}. Ergotropy itself is due to Allahverdyan,
Balian, and Nieuwenhuizen~\cite{Allahverdyan2004}.

\emph{Quantum energy teleportation.} The protocol and its spin-chain and field-theory
versions are Hotta's~\cite{Hotta2008,Hotta2009,Hotta2011review}, including energy--entanglement
relations for the minimal model~\cite{Hotta2010}. Ref.~\cite{Ragula2025} reviews experiments
and applications. Wang and Yao proved a trade-off between energy and information
teleportation~\cite{Wang2024}. Haque derived an upper bound on the energy extractable in gapped lattice systems with a
unique ground state that decays exponentially with the distance~\cite{Haque2026}. Itoh \emph{et al.} derived
information-thermodynamic bounds~\cite{Itoh2026}. Yan, Zhang, and Wang studied QET on mixed
steady states prepared by environments~\cite{Yan2026}. Xie, Sajjan, and Kais proposed storing
the extracted energy in a quantum register~\cite{Xie2024}, which is Bob's device made
quantum.

\emph{Closest work.} Xie, Sajjan, and Kais~\cite{Xie2025} relax three conventional constraints
on QET: a ground state, a measurement that commutes with the interaction, and entanglement.
They introduce Bob's local effective Hamiltonian, which is our $H^{(a)}$. In their protocol
Alice measures $\sigma_x^A$ while the coupling contains $\sigma_y^A\sigma_y^B$. That
measurement violates Eq.~\eqref{eq:ens}, since $\Abar^\dagger(\sigma_y^A)=0$, so it is not free
here, although on their initial state it injects no energy into the coupling. Their bond
operators $\sigma_x^A$ and $\sigma_y^A$ do not commute, so the model lies outside the scope of
Proposition~\ref{prop:closed} and Theorem~\ref{thm:freeclass}. Condition~\eqref{eq:ens} is
state-independent. For commuting bonds, Theorem~\ref{thm:freeclass} shows it is exactly the
commutation constraint whenever the bond spectrum is in convex position. Wu \emph{et al.}~\cite{Wu2024}
find strongly locally passive pure states of the minimal model by a correlation-matrix
analysis, with Bob's yield exceeding Alice's injection for some non-ground states. Fan
\emph{et al.}~\cite{Fan2024steering} analyse QET with the correlation matrix, identify
quantum steering as its essence, and define ``QET passive'' states. Those states are the
analogue of our free states; their class of allowed operations may differ. Fan
\emph{et al.}~\cite{Fan2024resources} study how quantum resources affect the efficiency of QET
and give necessary and sufficient conditions for the minimal model. These correlation-matrix
analyses may contain the two-qubit cases of Propositions~\ref{prop:closed}
and~\ref{prop:salvia}. We claim neither: Proposition~\ref{prop:salvia} is Salvia
\emph{et al.}'s, and in the minimal model Proposition~\ref{prop:closed} gives Hotta's value.
As far as we know, none of these works formulates a class of free operations with a ledger
monotone, characterizes the free class, or addresses two-way communication.

% ---------------------------------------------------------------------------------------
\section{Discussion and limitations}
\label{sec:discussion}

The theory's central object is a linear functional. For commuting bonds in convex
position, $\Wto(\rho)=\Tr[W_{\rm op}\rho]$, and everything Bob can receive by free means,
with one-way or two-way communication, is the expectation of one operator. The limitations
are the following.
\begin{itemize}\setlength{\itemsep}{1pt}\setlength{\parskip}{0pt}
\item Proposition~\ref{prop:closed} and Theorems~\ref{thm:freeclass} and~\ref{thm:linear}
assume commuting bond operators. Non-commuting bonds admit no finest free measurement, and we
do not treat them.
\item Off convex position we exhibit one free instrument outside the commuting class, not the
optimal one. Whether two-way communication helps there is open.
\item The theory credits Bob's side only. Alice's apparatus pays $E_A$, and crediting a second
device buys nothing net (Table~\ref{tab:twoway}).
\item For a single-qubit Bob, Eq.~\eqref{eq:salvia} is exact for unitaries, while the channel
optimum is a semidefinite program~\cite{Alhambra2019}. For a multi-qubit Bob we have no closed
form.
\item The entropic bounds are valid but loose by one to two orders of magnitude on the chains.
\item The numerics are finite-size: at most 12 sites and 4 modes. Nothing here is a
statement about a thermodynamic limit or a quantum field.
\end{itemize}

\begin{acknowledgments}
AI assistants, including Anthropic's Claude models, were used extensively in this work: to
write and test the analysis code, to run the numerical checks, to search the literature, and
to draft this manuscript. The author directed the work and takes full responsibility for its
content.
\end{acknowledgments}

% ---------------------------------------------------------------------------------------
\appendix

\section{Proofs}
\label{app:proofs}

\emph{Theorem~\ref{thm:ledger}(a).} Fix $\epsilon>0$. For each $\mu$ choose a free
instrument $\{K'_{\nu|\mu}\}$ and channels $\{\E_{\mu\nu}\}$ that deliver at least
$\Wto(\rho_\mu)-\epsilon$ from $\rho_\mu$. The composed instrument $\{K'_{\nu|\mu}K_\mu\}$ is
free, and with the channels $\E_{\mu\nu}$ it delivers
$\sum_\mu p_\mu[\Wto(\rho_\mu)-\epsilon]$ from $\rho$. Hence
$\Wto(\rho)\ge\sum_\mu p_\mu\Wto(\rho_\mu)-\epsilon$ for every $\epsilon$.

\emph{Theorem~\ref{thm:ledger}(b).} Let $\rho'=(\id\otimes\E)(\rho)$, fix $\epsilon>0$, and
choose $(\{K_\mu\},\{\E_\mu\})$ delivering at least $\Wto(\rho')-\epsilon$ from $\rho'$. Run
the same instrument on $\rho$, with Bob's channels $\E_\mu\circ\E$. Because $K_\mu$ acts on $A$
and $\E$ on $B$, $(K_\mu\otimes I)\rho'(K_\mu\otimes I)^\dagger=(\id\otimes\E)(\tilde\rho_\mu)$,
so the final energies of the two runs coincide. By Eq.~\eqref{eq:ens} their initial energies
are $\Tr[\HBl\rho]$ and $\Tr[\HBl\rho']$, which differ by $w_B(\E,\rho)$. Hence
$\Wto(\rho)\ge\Wto(\rho')-\epsilon+w_B(\E,\rho)$. The same arguments apply verbatim to $\WCK$.

\emph{Corollary~\ref{cor:ledger}.} Bob's steps obey Eq.~\eqref{eq:mono-b} and raise $D$ by
$w_B$. Alice's steps obey Eq.~\eqref{eq:mono-a} and leave $D$ unchanged. Communication
changes nothing. Mixing and discarding flags cannot increase $\Wto$, by convexity. Since
$\Wto\ge0$, the final $D$ is at most $\Wto(\rho)$.

\emph{Theorem~\ref{thm:freeclass}.} ($\Rightarrow$) Let every $x(b)$ be extreme, and let $v$
be a unit vector in the range of $\Pi_b$. By Eq.~\eqref{eq:blocks}, $p_a:=\bra vM_{ab}\ket v$
is a probability distribution with $\sum_ap_ax(a)=x(b)$. Suppose $p_b<1$. Then
$x(b)=\sum_{a\ne b}q_ax(a)$ with $q_a=p_a/(1-p_b)$, and taking any $a_0$ with $q_{a_0}>0$
writes $x(b)$ as a combination of $x(a_0)\ne x(b)$ and a point of the hull. That is either
trivial, which would give $x(a_0)=x(b)$, or contradicts extremality. So $p_b=1$ and $p_a=0$
for $a\ne b$. Since $M_{ab}\ge0$ is supported in the range of $\Pi_b$ and $v$ was arbitrary,
$M_{ab}=0$, i.e., $\sum_\mu\|\Pi_aK_\mu\Pi_b\psi\|^2=0$ for all $\psi$. Every $K_\mu$ is
then block diagonal and commutes with every $\Pi_a$, hence with every $X_j$. ($\Leftarrow$)
If some $x(b)$ is not extreme, it lies in $\mathrm{conv}\{x(a)\}_{a\ne b}$, because the extreme
points of the hull of a finite set belong to the set. The instrument~\eqref{eq:witness} is
then complete, satisfies Eq.~\eqref{eq:ens}, and has $\Pi_aK_a\Pi_b\ne0$.

\emph{Proposition~\ref{prop:surplus}.} Within the proposition's restriction the potential is
$\WCK$, so $Q=\WCK-\Wloc$ (in the minimal model $\WCK=\Wto$). Write
$m(\rho):=\min_\E\Tr[\HBl(\id\otimes\E)(\rho)]$, so that $\Wloc(\rho)=\Tr[\HBl\rho]-m(\rho)$,
and $\Lambda(\rho):=\sum_ap_a\lmin(H^{(a)})$, so that $\WCK(\rho)=\Tr[\HBl\rho]-\Lambda(\rho)$
by Eq.~\eqref{eq:closed}. Hence $Q=m(\rho)-\Lambda(\rho)$.
(a) A channel $\E'$ on $B$ leaves the $p_a$ and hence $\Lambda$ unchanged, and
$m(\E'(\rho))\ge m(\rho)$ because $\E\circ\E'$ is a channel. So $Q$ does not decrease. For a
commuting-Kraus instrument, $\sum_\mu K_\mu^\dagger\Pi_aK_\mu=\Pi_a$ gives
$\sum_\mu p_\mu\Lambda(\rho_\mu)=\Lambda(\rho)$. Let $\E^\star$ attain $m(\rho)$. Then
$\sum_\mu p_\mu m(\rho_\mu)\le\sum_\mu\Tr[\HBl(\id\otimes\E^\star)(\tilde\rho_\mu)]=m(\rho)$
by Eq.~\eqref{eq:ens}. So $Q$ does not increase on average.
(b) For a product state, $m=\lmin(\bar H)$ with the mean field
$\bar H=h\sigma_z+2k\langle\sigma_x^A\rangle\sigma_x$. For $\rho_1$: $\Wto=\Wloc=h+\sqrt{h^2+4k^2}$.
For $\rho_2$: $\Wto=h+\sqrt{h^2+4k^2}$, since Alice's $\sigma_x$ measurement restores the
field $\pm2k\sigma_x$ on each branch, while $\Wloc=2h$.
(c) For $\rho_3$, Bob can measure $\sigma_z$ and prepare the ground state of $H^{(\pm)}$,
reaching $m=\Lambda=-\sqrt{h^2+4k^2}$, so $Q=0$. $\Tr[\HBl\rho_3]=0$ and
$\Tr[\HBl\rho_5]=-h$, so the reset delivers $h$. For the product state $\rho_5$,
$\Wto=\sqrt{h^2+4k^2}-h$ and $\Wloc=0$.

\emph{Proposition~\ref{prop:entropic}.} (a) $U\tau U^\dagger$ has largest eigenvalue
$e^{-\Hmin(\tau)}$, so it is feasible in the definition of $\kappa$. The minimum is attained
on states diagonal in the eigenbasis of $H$ by filling the lowest levels to the cap
$e^{-s}$. (b) Let $Y_U:=H_R-(U^\dagger\otimes I)H_R(U\otimes I)$. Passivity gives
$\Tr[Y_U\rho_R]\le0$ for every $U$. So
$E_B=\sum_\mu p_\mu\max_U\Tr[Y_U\tau_\mu]\le\sum_\mu p_\mu\max_U\Tr[Y_U(\tau_\mu-\rho_R)]$.
For traceless $\delta$, $|\Tr[Y\delta]|\le\tfrac12\spread(Y)\|\delta\|_1$, so this is at most
$c\sum_\mu p_\mu\|\tau_\mu-\rho_R\|_1=\tfrac c2\|\tau_+-\tau_-\|_1$. The Helstrom formula
$e^{-\Hmin(X|R)}=\tfrac12+\tfrac14\|\tau_+-\tau_-\|_1$~\cite{KRS2009} gives Eq.~\eqref{eq:guess}.

\section{Code and data}
\label{app:code}

Every number in this paper comes from archived data produced by one script. The script is
\url{papers/resource-theory-vacuum-manipulation/notebook.py} (SHA-256
\texttt{ddd18909\ldots e68f}), run on 5 September 2026 on a tree whose full test suite passed.
The candidate's own tests are in \url{tests/test_resource_theory.py} (44 tests). The
figure tables are written from the archive by \url{export_figdata.py} next to this
manuscript. Three statements were checked separately while this draft was written: the
certified gap at $L=6$ (Sec.~\ref{sec:single}), the identity of Eq.~\eqref{eq:EBmax} with
Hotta's formula, and Proposition~\ref{prop:surplus}(c). The code and data are available at
\url{https://github.com/mruckman1/vacuum_energy}. Its numbering differs from the paper's, as
follows.

\begin{table}[h]
\caption{Correspondence between this paper and the module docstring of
\texttt{vacuum/qet/resource.py}.}
\label{tab:map}
\begin{ruledtabular}
\begin{tabular}{ll}
This paper & Code \\
\hline
Eq.~\eqref{eq:ens}, energy non-signalling & Eq.~(2)\\
Eq.~\eqref{eq:W}, the potential & Eq.~(3)\\
Proposition~\ref{prop:closed} & Proposition A, Eq.~(4)\\
Theorem~\ref{thm:ledger} & Theorem 1, Eqs.~(5a)--(5b)\\
Theorem~\ref{thm:freeclass}, Eqs.~\eqref{eq:blocks}--\eqref{eq:witness} & Theorem 2, Eqs.~(7)--(8)\\
Eq.~\eqref{eq:sdp} & Eq.~(20)\\
Theorem~\ref{thm:linear} & Theorem 4, Eqs.~(10)--(11)\\
Corollary~\ref{cor:twoway} & Eq.~(12)\\
Proposition~\ref{prop:surplus} & surplus section\\
Proposition~\ref{prop:salvia} & Proposition B, Eq.~(15)\\
Proposition~\ref{prop:entropic}(a) & chain (ii)\\
Proposition~\ref{prop:entropic}(b) & Proposition C, Eq.~(17)\\
Eq.~\eqref{eq:gauss} & Gaussian mirror\\
\end{tabular}
\end{ruledtabular}
\end{table}

\bibliography{refs}

\end{document}
